% GENERATED FILE. Edit canonical lessons, then run npm run generate:books.
% canonical-source-hash: fnv1a64:9d418130ddf648c9
% canonical-lessons: JA-C283-happyou-suru, JA-C283-houkoku-suru, JA-C283-chousa-suru, JA-C283-jouhou, JA-C283-jijitsu

\chapter{The Announcement}
\label{ch:ja-the-announcement}
\hlchaptermodality{283}

\begin{chapteropening}
\textbf{By the end of this chapter:} \emph{I can say what the authorities announced and looked into.}

Complete the last lesson of chapter 283: I can say what the authorities announced and looked into.
\end{chapteropening}

\section[\texorpdfstring{\ja{はっぴょうする} (happyō suru)}{はっぴょうする (happyō suru)}]{\ja{はっぴょうする} (happyō suru) — to announce; to make public}

\label{lesson:JA-C283-happyou-suru}



\begin{quote}
Before the new one: say the Japanese for to fight, then the Japanese for the government.
\end{quote}

\subsection*{You'll want to know: \ja{はっぴょうする}}
\textbf{\ja{はっぴょうする}} — \emph{happyō suru} — \textquotedblleft{}to announce; to make public\textquotedblright{}.

\begin{quote}
\textbf{\ja{せいふは} \ja{あした} \ja{はっぴょうする} \ja{つもりです。}} — \emph{seifu wa ashita happyō suru tsumori desu} — \textquotedblleft{}The government means to make an announcement tomorrow.\textquotedblright{}
\end{quote}

Polite \textbf{\ja{はっぴょうします}}, \emph{happyōshimasu}; a \ja{する} verb. Past (polite): \textbf{\ja{せいふが} \ja{はっぴょうしました。}}, \emph{seifu ga happyō shimashita}.

\begin{cousinweb}
\textbf{\ja{はっぴょうする}} (\emph{happyō suru}) is Sino-Japanese (kango): two Chinese characters read with their on'yomi, readings borrowed from Chinese. They are \emph{hatsu}, \textquotedblleft{}publish\textquotedblright{} (Middle Chinese \emph{pjot}) and \emph{hyō}, \textquotedblleft{}surface\textquotedblright{} (\emph{pjewX}) (KANJIDIC2; Middle Chinese in Baxter's transcription). Together they are read \emph{happyō}: \emph{hatsu} shortens to \emph{hap-} and \emph{hyō} becomes \emph{pyō}. The native verb \emph{suru}, \textquotedblleft{}to do\textquotedblright{}, makes it a verb.
\end{cousinweb}

\begin{grammarlens}[title={What you've built}]
One more word for telling what happened.
\end{grammarlens}

\subsection*{Your turn}
\begin{itemize}
\raggedright
  \item \emph{Say it:} \emph{happyō suru}
  \item \emph{Say it:} \emph{happyō suru}, once more
  \item \emph{Say it:} \emph{seifu wa ashita happyō suru tsumori desu} — The government means to make an announcement tomorrow.
  \item \emph{Recall:} say the Japanese for to fight, then the Japanese for the government, then say \emph{happyō suru} again
\end{itemize}

\begin{tcolorbox}[breakable,colback=teal!4,colframe=teal!35!black,title={Before you move on}]
What does \ja{はっぴょうする} mean? (\textquotedblleft{}To announce; to make public\textquotedblright{}.) Say it once more.
\end{tcolorbox}
\section[\texorpdfstring{\ja{ほうこくする} (hōkoku suru)}{ほうこくする (hōkoku suru)}]{\ja{ほうこくする} (hōkoku suru) — to report; to give an account}

\label{lesson:JA-C283-houkoku-suru}



\begin{quote}
Before the new one: say the Japanese for the government, then the Japanese for to announce.
\end{quote}

\subsection*{You'll want to know: \ja{ほうこくする}}
\textbf{\ja{ほうこくする}} — \emph{hōkoku suru} — \textquotedblleft{}to report; to give an account\textquotedblright{}.

\begin{quote}
\textbf{\ja{おまわりさんが} \ja{じこを} \ja{ほうこくしました。}} — \emph{omawarisan ga jiko o hōkoku shimashita} — \textquotedblleft{}The officer reported the accident.\textquotedblright{}
\end{quote}

Polite \textbf{\ja{ほうこくします}}, \emph{hōkokushimasu}; a \ja{する} verb. Past (polite): \textbf{\ja{きしゃが} \ja{ほうこくしました。}}, \emph{kisha ga hōkoku shimashita}.

\begin{cousinweb}
\textbf{\ja{ほうこくする}} (\emph{hōkoku suru}) is Sino-Japanese (kango): two Chinese characters read with their on'yomi, readings borrowed from Chinese. They are \emph{hō}, \textquotedblleft{}report\textquotedblright{} (Middle Chinese \emph{pawH}) and \emph{koku}, \textquotedblleft{}inform\textquotedblright{} (\emph{kowk}) (KANJIDIC2; Middle Chinese in Baxter's transcription). The native verb \emph{suru}, \textquotedblleft{}to do\textquotedblright{}, makes it a verb.
\end{cousinweb}

\begin{grammarlens}[title={What you've built}]
One more word for telling what happened.
\end{grammarlens}

\subsection*{Your turn}
\begin{itemize}
\raggedright
  \item \emph{Say it:} \emph{hōkoku suru}
  \item \emph{Say it:} \emph{hōkoku suru}, once more
  \item \emph{Say it:} \emph{omawarisan ga jiko o hōkoku shimashita} — The officer reported the accident.
  \item \emph{Recall:} say the Japanese for the government, then the Japanese for to announce, then say \emph{hōkoku suru} again
\end{itemize}

\begin{tcolorbox}[breakable,colback=teal!4,colframe=teal!35!black,title={Before you move on}]
What does \ja{ほうこくする} mean? (\textquotedblleft{}To report; to give an account\textquotedblright{}.) Say it once more.
\end{tcolorbox}
\section[\texorpdfstring{\ja{ちょうさする} (chōsa suru)}{ちょうさする (chōsa suru)}]{\ja{ちょうさする} (chōsa suru) — to investigate; to survey}

\label{lesson:JA-C283-chousa-suru}



\begin{quote}
Before the new one: say the Japanese for to announce, then the Japanese for to report.
\end{quote}

\subsection*{You'll want to know: \ja{ちょうさする}}
\textbf{\ja{ちょうさする}} — \emph{chōsa suru} — \textquotedblleft{}to investigate; to survey\textquotedblright{}.

\begin{quote}
\textbf{\ja{けいさつが} \ja{じこの} \ja{ばしょを} \ja{ちょうさしています。}} — \emph{keisatsu ga jiko no basho o chōsa shite imasu} — \textquotedblleft{}The police are examining the scene of the accident.\textquotedblright{}
\end{quote}

Polite \textbf{\ja{ちょうさします}}, \emph{chōsashimasu}; a \ja{する} verb. Past (polite): \textbf{\ja{けいさつが} \ja{ちょうさしました。}}, \emph{keisatsu ga chōsa shimashita}.

\begin{cousinweb}
\textbf{\ja{ちょうさする}} (\emph{chōsa suru}) is Sino-Japanese (kango): two Chinese characters read with their on'yomi, readings borrowed from Chinese. They are \emph{chō}, \textquotedblleft{}prepare\textquotedblright{} (Middle Chinese \emph{dewH}) and \emph{sa}, \textquotedblleft{}investigate\textquotedblright{} (KANJIDIC2; Middle Chinese in Baxter's transcription). The native verb \emph{suru}, \textquotedblleft{}to do\textquotedblright{}, makes it a verb.
\end{cousinweb}

\begin{grammarlens}[title={What you've built}]
One more word for telling what happened.
\end{grammarlens}

\subsection*{Your turn}
\begin{itemize}
\raggedright
  \item \emph{Say it:} \emph{chōsa suru}
  \item \emph{Say it:} \emph{chōsa suru}, once more
  \item \emph{Say it:} \emph{keisatsu ga jiko no basho o chōsa shite imasu} — The police are examining the scene of the accident.
  \item \emph{Recall:} say the Japanese for to announce, then the Japanese for to report, then say \emph{chōsa suru} again
\end{itemize}

\begin{tcolorbox}[breakable,colback=teal!4,colframe=teal!35!black,title={Before you move on}]
What does \ja{ちょうさする} mean? (\textquotedblleft{}To investigate; to survey\textquotedblright{}.) Say it once more.
\end{tcolorbox}
\section[\texorpdfstring{\ja{じょうほう} (jōhō)}{じょうほう (jōhō)}]{\ja{じょうほう} (jōhō) — information; news}

\label{lesson:JA-C283-jouhou}



\begin{quote}
Before the new one: say the Japanese for to report, then the Japanese for to investigate.
\end{quote}

\subsection*{You'll want to know: \ja{じょうほう}}
\textbf{\ja{じょうほう}} — \emph{jōhō} — \textquotedblleft{}information; news\textquotedblright{}.

\begin{quote}
\textbf{\ja{じょうほうが} \ja{すくないです。}} — \emph{jōhō ga sukunai desu} — \textquotedblleft{}There is little to go on.\textquotedblright{}
\end{quote}

\begin{cousinweb}
\textbf{\ja{じょうほう}} (\emph{jōhō}) is Sino-Japanese (kango): two Chinese characters read with their on'yomi, readings borrowed from Chinese. They are \emph{jō}, \textquotedblleft{}circumstances\textquotedblright{} (Middle Chinese \emph{dzjeng}) and \emph{hō}, \textquotedblleft{}report\textquotedblright{} (\emph{pawH}) (KANJIDIC2; Middle Chinese in Baxter's transcription).
\end{cousinweb}

\begin{grammarlens}[title={What you've built}]
One more word for telling what happened.
\end{grammarlens}

\subsection*{Your turn}
\begin{itemize}
\raggedright
  \item \emph{Say it:} \emph{jōhō}
  \item \emph{Say it:} \emph{jōhō}, once more
  \item \emph{Say it:} \emph{jōhō ga sukunai desu} — There is little to go on.
  \item \emph{Recall:} say the Japanese for to report, then the Japanese for to investigate, then say \emph{jōhō} again
\end{itemize}

\begin{tcolorbox}[breakable,colback=teal!4,colframe=teal!35!black,title={Before you move on}]
What does \ja{じょうほう} mean? (\textquotedblleft{}Information; news\textquotedblright{}.) Say it once more.
\end{tcolorbox}
\section[\texorpdfstring{\ja{じじつ} (jijitsu)}{じじつ (jijitsu)}]{\ja{じじつ} (jijitsu) — a fact; the truth of the matter}

\label{lesson:JA-C283-jijitsu}



\begin{quote}
Before the new one: say the Japanese for to investigate, then the Japanese for information.
\end{quote}

\subsection*{You'll want to know: \ja{じじつ}}
\textbf{\ja{じじつ}} — \emph{jijitsu} — \textquotedblleft{}a fact; the truth of the matter\textquotedblright{}.

\begin{quote}
\textbf{\ja{けいさつは} \ja{じじつを} \ja{しらべています。}} — \emph{keisatsu wa jijitsu o shirabete imasu} — \textquotedblleft{}The police are looking into the facts.\textquotedblright{}
\end{quote}

\begin{cousinweb}
\textbf{\ja{じじつ}} (\emph{jijitsu}) is Sino-Japanese (kango): two Chinese characters read with their on'yomi, readings borrowed from Chinese. They are \emph{ji}, \textquotedblleft{}matter\textquotedblright{} (Middle Chinese \emph{dzriH}) and \emph{jitsu}, \textquotedblleft{}reality\textquotedblright{} (\emph{zyit}) (KANJIDIC2; Middle Chinese in Baxter's transcription).
\end{cousinweb}

\begin{grammarlens}[title={What you've built}]
One more word for telling what happened.
\end{grammarlens}

\subsection*{Your turn}
\begin{itemize}
\raggedright
  \item \emph{Say it:} \emph{jijitsu}
  \item \emph{Say it:} \emph{jijitsu}, once more
  \item \emph{Say it:} \emph{keisatsu wa jijitsu o shirabete imasu} — The police are looking into the facts.
  \item \emph{Recall:} say the Japanese for to investigate, then the Japanese for information, then say \emph{jijitsu} again
\end{itemize}

\begin{tcolorbox}[breakable,colback=teal!4,colframe=teal!35!black,title={Before you move on}]
What does \ja{じじつ} mean? (\textquotedblleft{}A fact; the truth of the matter\textquotedblright{}.) Say it once more.
\end{tcolorbox}
